\documentclass[11pt]{article}
\usepackage[margin=1in]{geometry}
\usepackage{amsmath,amssymb,amsthm}
\usepackage[htt]{hyphenat}
\usepackage[hidelinks]{hyperref}
\usepackage{booktabs}
\usepackage{listings}
\hyphenation{ap-prox-i-mant tele-scop-ing mul-ti-mod-al-ity in-jec-tiv-ity}
\sloppy

\newtheorem{theorem}{Theorem}[section]
\newtheorem{proposition}[theorem]{Proposition}
\newtheorem{lemma}[theorem]{Lemma}
\newtheorem{corollary}[theorem]{Corollary}
\theoremstyle{definition}
\newtheorem{assumption}[theorem]{Assumption}
\newtheorem{definition}[theorem]{Definition}
\newtheorem{conjecture}[theorem]{Conjecture}
\theoremstyle{remark}
\newtheorem{remark}[theorem]{Remark}

\newcommand{\E}{\mathbb{E}}
\newcommand{\Var}{\operatorname{Var}}
\newcommand{\R}{\mathbb{R}}
\DeclareMathOperator{\dH}{d_{H}}

\title{Telescoping Bayesian Inference with Approximate Forward Operators:\\
Corrected Convergence Theory, Multimodality as Loss of Injectivity,\\ and the Cost of Heavy Tails}
\author{Josh Bald\\ \small Independent Researcher\\ \small jpbald93@gmail.com%
\thanks{Manuscript source, the verification driver, and the two independent audit reports are
available at \url{https://github.com/jpbald93/bayesian-telescoping-v2}.}}
\date{September 2026 \quad (V2)}

\begin{document}
\maketitle

\begin{abstract}
Multilevel and telescoping methods for Bayesian inverse problems are well developed under ideal
conditions. This paper is a \emph{corrected} revision of an earlier manuscript that claimed three
results beyond the standard theory. Two of the three claims do not survive scrutiny; the third is
provable only in a restricted form, which we give. Specifically, (i) the multilevel regime
thresholds are $\gamma=\beta+2$, not $\gamma=\beta+1$, and the sub-optimal exponent is
$-2-(\beta+2-\gamma)/k$ (the bias is $O(L^{-k})$, so $L\asymp\varepsilon^{-1/k}$); (ii) discretization-induced bimodality \emph{cannot} arise from a
linear forward operator with a Gaussian prior --- the approximate posterior is then exactly Gaussian
at every level --- and for the Runge--Kutta~4 surrogate of $\dot y=\lambda y$ it is a
\emph{coarse-resolution transient} caused by non-injectivity of the numerical map, which disappears
once the fold $\lambda=z^\star n/T$ leaves the prior's effective support, rather than persisting at
all levels; and
(iii) the claimed heavy-tail penalty $\rho(\alpha)=(\alpha-2)/(\alpha-1)$ is neither proved nor
supported by the earlier paper's own experiment, which we reproduce and which in fact shows the full
rate $-(k+1)$ for every tail index tested. We state precisely what is proved, what is verified
numerically, and what remains open, and we ship a verification driver that regenerates every numeral
in this paper.
\end{abstract}

\section{Introduction}

\subsection{What this paper is}
This is the corrected second version of a manuscript (hereafter \textbf{V1}) on telescoping
Bayesian inference with approximate forward operators. V1's classical material (telescoping
posterior decomposition, expectation convergence, variance decay, credible-interval convergence) is
essentially correct and is retained, with repairs. V1's three advertised novel contributions,
however, do not hold up:
\begin{enumerate}
\item its multilevel complexity regimes are off by one (\S\ref{sec:mlmc});
\item its theory of discretization-induced bimodality rests on an invalid Hessian argument and on a
      characterisation contradicted by its own example, and its empirical ``confirmation'' is a
      simulation whose code hard-codes the bimodality it claims to explain (\S\ref{sec:bimodal});
\item its heavy-tail penalty formula is asserted rather than derived, is attributed to a
      mis-cited reference, and is not exhibited by its own experiment design (\S\ref{sec:tails}).
\end{enumerate}
We therefore \emph{retract} V1's Theorems 9.3 and 10.1 and its Theorem 8.4 as stated, and replace
them with the correct statements below. Two of the replacement statements are theorems proved here;
the third is explicitly posed as a conjecture with numerical evidence against V1's formula.

\subsection{Status of the claims}
\begin{center}
\begin{tabular}{lll}
\toprule
Claim & Status in V2 & Where \\
\midrule
Telescoping decomposition & correct (elementary identity) & \S\ref{sec:frame} \\
Expectation convergence $O(n^{-(k+1)})$ & correct (unchanged) & Thm~\ref{thm:conv} \\
Variance-difference decay & correct (unchanged) & Thm~\ref{thm:var} \\
MLMC regimes & \textbf{corrected} (thresholds $\gamma=\beta+2$) & Thm~\ref{thm:mlmc} \\
Bimodality from linear $F$ & \textbf{impossible} (new, proved) & Thm~\ref{thm:obstruction} \\
Bimodality from RK4 & \textbf{transient}, not persistent (new) & \S\ref{sec:rk4} \\
Heavy-tail penalty $\rho(\alpha)$ & \textbf{not proved}; not observed numerically & \S\ref{sec:tails} \\
\bottomrule
\end{tabular}
\end{center}

\subsection{Notation}
$X$ is a separable Banach space (parameter), $Y$ a separable Hilbert space (observations),
$y=F(x)+\eta$ with $\eta\sim N(0,\Gamma)$. $\mu_0$ is a prior, $\mu^y$ the exact posterior, and
$\mu_n^y$ the posterior induced by an approximate forward operator $F_n$. $k$ is the telescoping
order: $\|F_{n+1}(x)-F_n(x)\|\le C(1+\|x\|^p)/n^{k+1}$.

\section{Preliminaries}\label{sec:prelim}
We use the Hellinger distance
$\dH(\mu,\nu)^2=\int(\sqrt{d\mu/d\lambda}-\sqrt{d\nu/d\lambda})^2\,d\lambda$, which
satisfies $\dH\in[0,\sqrt2]$ and $\dH=\sqrt{2(1-\rho)}$ with the Bhattacharyya coefficient $\rho$.
(With the alternative normalisation $\tfrac12\int(\cdot)^2$ one has $\dH\in[0,1]$ and
$\dH=\sqrt{1-\rho}$; V1 mixed the two conventions.)
For a Gaussian prior satisfying Fernique's theorem there is $a>0$ with
$\int\exp(a\|x\|^2)\,\mu_0(dx)<\infty$; we shall need the stronger, and separate, \emph{uniform}
exponential-moment condition
\begin{equation}\label{eq:unif}
  \sup_{n\ge1}\int_X \exp\bigl(a\|x\|^{2p}\bigr)\,\mu_n^y(dx)<\infty
\end{equation}
with $p$ the polynomial-growth exponent. The potential is
$\Phi_n(x;y)=\tfrac12\|\Gamma^{-1/2}(y-F_n(x))\|_Y^2$ and
$d\mu_n^y/d\mu_0 = Z_n(y)^{-1}\exp(-\Phi_n(\cdot;y))$.

We record the elementary facts we shall use, with proofs, because V1 got one of them wrong.

\begin{lemma}[Telescoping tail; corrected]\label{lem:tail}
Let $\alpha>1$. Then
\[
  \sum_{n\ge N} n^{-\alpha} \;\le\; N^{-\alpha} + \int_N^\infty t^{-\alpha}\,dt
  \;=\; N^{-\alpha} + \frac{N^{1-\alpha}}{\alpha-1} \;\le\; \frac{C_\alpha}{N^{\alpha-1}},
\]
and in particular $\sum_{n\ge N} n^{-\alpha}=\Theta(N^{1-\alpha})$ (the \emph{tail}; note that
$\sum_{n<N}n^{-\alpha}=\Theta(1)$ for $\alpha>1$).
\end{lemma}
\begin{proof}
For decreasing $t\mapsto t^{-\alpha}$ we have $n^{-\alpha}\le\int_{n-1}^{n}t^{-\alpha}dt$ for
$n\ge2$, so $\sum_{n\ge N+1}n^{-\alpha}\le\int_N^\infty t^{-\alpha}dt$, and adding the $n=N$ term
gives the first bound; the second is immediate for $N\ge1$ with $C_\alpha=1+1/(\alpha-1)$.
\end{proof}
\begin{remark}
V1 bounded $\int_{N-1}^\infty t^{-\alpha}dt=1/((\alpha-1)(N-1)^{\alpha-1})$ and then wrote
``$\le1/((\alpha-1)N^{\alpha-1})$'', which is the wrong direction: the sum exceeds
$1/((\alpha-1)N^{\alpha-1})$ (e.g.\ $\alpha=2,N=10$: $0.1052>0.1000$). The driver asserts both
sides. The conclusion $\Theta(N^{-(\alpha-1)})$ is unaffected; only the displayed step was wrong.
\end{remark}

\section{Telescoping posteriors}\label{sec:frame}

\begin{proposition}[Increment identity; correct]\label{prop:inc}
$\displaystyle \Phi_{n+1}(x;y)-\Phi_n(x;y)
   =-\langle\Gamma^{-1}(y-F_n(x)),\,F_{n+1}(x)-F_n(x)\rangle
    +\tfrac12\|\Gamma^{-1/2}(F_{n+1}(x)-F_n(x))\|^2.$
\end{proposition}
\begin{proof}
Write $a=y-F_n(x)$, $d=F_{n+1}(x)-F_n(x)$, so $y-F_{n+1}(x)=a-d$. Then
$\Phi_{n+1}=\tfrac12\langle\Gamma^{-1}(a-d),a-d\rangle
=\tfrac12\langle\Gamma^{-1}a,a\rangle-\langle\Gamma^{-1}a,d\rangle+\tfrac12\langle\Gamma^{-1}d,d\rangle$.
\end{proof}

\begin{corollary}\label{cor:incbound}
Assume $\|F_n(x)\|_Y\le C(1+\|x\|_X^p)$ uniformly in $n$. Then
$|\Phi_{n+1}(x;y)-\Phi_n(x;y)|\le C(1+\|x\|_X^{2p})(1+\|y\|_Y)n^{-(k+1)}+C(1+\|x\|_X^{2p})n^{-2(k+1)}$.
\end{corollary}

A standard dominated-convergence argument gives the following (V1's Theorem 5.2), which we do not
repeat in full.

\begin{theorem}[Posterior convergence of expectations]\label{thm:conv}
Under the assumed growth and \eqref{eq:unif}, for every bounded Lipschitz $\varphi$ there is $C$
with $|\E_{\mu^y_{n+1}}[\varphi]-\E_{\mu^y_n}[\varphi]|\le Cn^{-(k+1)}$, and hence
$|\E_{\mu^y}[\varphi]-\E_{\mu^y_N}[\varphi]|\le C_1(kN^k)^{-1}$.
\end{theorem}

\begin{theorem}[Variance-difference decay]\label{thm:var}
Under the hypotheses of Theorem~\ref{thm:conv}, $|\Var_{\mu^y_{n+1}}(\varphi)-\Var_{\mu^y_n}(\varphi)|=O(n^{-(k+1)})$.
\end{theorem}

\begin{proposition}[Credible intervals; correct hypotheses]\label{prop:ci}
Let $\varphi$ be bounded with marginal densities $p_n$ under $\mu^y_n$, and let $F_n$ be the
distribution function of $\varphi$ under $\mu^y_n$, with $\alpha$-quantile $q_\alpha^{(n)}$.
Suppose (i) $\inf_n\inf_{|s-q_\alpha|<\delta}p_n(s)\ge c>0$ for some $\delta>0$, and (ii)
$|F_n(q_\alpha)-\alpha|=O(\delta_n)$, where $q_\alpha=F^{-1}(\alpha)$. Then
$|q_\alpha^{(n)}-q_\alpha|=O(\delta_n)$.
\end{proposition}
\begin{proof}
By the mean value theorem applied to $F_n$,
$|q_\alpha^{(n)}-q_\alpha|=|F_n(q_\alpha)-\alpha|/p_n(\xi_n)$ for some $\xi_n$ between the two
quantiles. The numerator is $O(\delta_n)$ by (ii); the denominator is $\ge c$ by (i) once
$|q_\alpha^{(n)}-q_\alpha|<\delta$.
\end{proof}
\begin{remark}
Hypothesis (ii) is \emph{not} implied by convergence of the first two moments: two log-concave
densities with identical mean and variance may have different quantiles. V1's Corollary 7.2 assumed
only moment convergence and is therefore false; what is needed --- and what V1's own Lemma A.1 in
fact used --- is convergence of the distribution function at the quantile. Note also that (i) is a
hypothesis \emph{about the approximate posteriors}; it is automatic when they are log-concave, which
\S\ref{sec:bimodal} shows holds for affine $F_n$ and Gaussian priors, but it can fail in general.
\end{remark}

\section{Multilevel complexity: corrected regimes}\label{sec:mlmc}
Let $V_n=\Var(\Delta_n)=O(n^{-\gamma})$ and $C_n=O(n^{\beta})$, and let the estimator use $L$ levels
with optimal allocation $N_n\propto\sqrt{V_n/C_n}$. The standard identity
$\mathrm{Cost}(\varepsilon)\asymp\varepsilon^{-2}\bigl(\sum_{n<L}\sqrt{V_nC_n}\bigr)^2$ with
$L\asymp\varepsilon^{-1/k}$ gives, since $\sqrt{V_nC_n}\asymp n^{(\beta-\gamma)/2}$,

\begin{theorem}[MLMC complexity; corrected]\label{thm:mlmc}
Write $s=(\beta-\gamma)/2$ and $L\asymp\varepsilon^{-1/k}$.
\begin{enumerate}
  \item If $\gamma>\beta+2$: $\ \sum_{n<L}n^{s}=O(1)$ and $\mathrm{Cost}(\varepsilon)=O(\varepsilon^{-2})$.
  \item If $\gamma=\beta+2$: $\ \sum_{n<L}n^{s}=\Theta(\log L)$ and
        $\mathrm{Cost}(\varepsilon)=O\bigl(\varepsilon^{-2}(\log\varepsilon)^2\bigr)$.
  \item If $\gamma<\beta+2$: $\ \sum_{n<L}n^{s}=\Theta\bigl(L^{(\beta+2-\gamma)/2}\bigr)$ and
        $\mathrm{Cost}(\varepsilon)=O\bigl(\varepsilon^{-2-(\beta+2-\gamma)/k}\bigr)$.
\end{enumerate}
\end{theorem}
\begin{proof}
$\sum_{n<L}n^{s}$ converges iff $s<-1$, i.e.\ $\gamma>\beta+2$; at $s=-1$ it is $\Theta(\log L)$;
for $s>-1$ it is $\Theta(L^{s+1})=\Theta(L^{(\beta+2-\gamma)/2})$. Substituting
$L\asymp\varepsilon^{-1/k}$ gives (iii).
\end{proof}
\begin{remark}[Cutoff: the bias is $O(L^{-k})$, not $O(L^{-(k+1)})$]
Theorem~\ref{thm:conv} bounds the \emph{increment} by $Cn^{-(k+1)}$; the level-$L$ bias is the tail
sum $\sum_{n\ge L}|\E_{\mu_{n+1}}[\varphi]-\E_{\mu_n}[\varphi]|=O(L^{-k})$ by Lemma~\ref{lem:tail}
with $\alpha=k+1$. Hence $L\asymp\varepsilon^{-1/k}$, not $\varepsilon^{-1/(k+1)}$. (The driver
verifies $\sum_{n\ge L}n^{-(k+1)}=\Theta(L^{-k})$.) V1 used $L\asymp\varepsilon^{-1/(k+1)}$, which is
inconsistent with its own bias corollary; we correct the cutoff here and correspondingly the
exponent in case (iii).
\end{remark}

\begin{remark}[What V1 said, and why it is off by one]
V1 asserted the threshold $\gamma>\beta+1$, the boundary case $\gamma=\beta+1\Rightarrow O(\log L)$,
and the sub-optimal exponent $-2-(\beta+1-\gamma)/(k+1)$. Its own derivation contains the slip
$(\beta-\gamma)/2+1=(\beta+1-\gamma)/2$, which is false: the left side is
$(\beta+2-\gamma)/2$ (e.g.\ $\beta=1,\gamma=3$: left $=0$, right $=-\tfrac12$). With V1's threshold,
the partial sums diverge: the driver reports $\sum_{n<10^6}n^{-3/4}\approx123$ at the V1 boundary,
i.e.\ V1's ``$O(1)$'' case contains divergent sums. The empirical derivation of $\gamma$ from $k$
(``$\gamma=2(k+1)$'') is in V1 and remains an assumption, not a theorem.
\end{remark}

\section{Discretization error and multimodality}\label{sec:bimodal}
V1's Theorem 9.3 claimed that ``oscillatory error'' (RK4) causes bimodality whereas ``sign-definite
error with $k\ge2$'' prevents it, and offered a Hessian argument. That argument is invalid: the
Hessian of $-\log(d\mu^y_n/d\mu_0)$ omits the prior curvature, and the middle term below is missing
from V1. Let $H_n$ be the Hessian of the negative logarithm of the \emph{density} of $\mu^y_n$ with
respect to Lebesgue measure (prior included):
\begin{equation}\label{eq:hess}
  H_n(x)=\nabla F_n(x)^\top\Gamma^{-1}\nabla F_n(x)
        -\bigl\langle\Gamma^{-1}(y-F_n(x)),\nabla^2F_n(x)\bigr\rangle
        +\text{Prior}^{-1},
\end{equation}
and V1 dropped the middle term, which is indefinite and $O(1)$ in $n$. Two of V1's claims also
contradict its own example: RK4 is a \emph{fourth}-order method (V1 labels it $k=1$), and the
``odd/even-order'' dichotomy is not the same as the ``$k<2$ / $k\ge2$'' dichotomy.

We replace Theorem 9.3 with two correct statements.

\subsection{An obstruction: linear forward operators cannot produce bimodality}\label{sec:lin}
\begin{theorem}[No bimodality for linear $F$; new]\label{thm:obstruction}
Let $F_n(x)=A_nx+b_n$ be affine and let the prior be Gaussian, $\mu_0=N(m_0,\Sigma_0)$, with
Gaussian noise. Then for every $n$ the approximate posterior $\mu_n^y$ is Gaussian, hence
log-concave and unimodal; in particular no choice of discretisation of a \emph{linear} forward
operator can induce a bimodal posterior at any level.
\end{theorem}
\begin{proof}
$\Gamma^{-1/2}(y-A_nx-b_n)$ is affine in $x$, so $-\Phi_n$ is a concave quadratic in $x$ and
$\exp(-\Phi_n)d\mu_0$ is (up to a constant) a Gaussian density; hence $\mu_n^y$ is Gaussian.
Gaussian densities are strictly log-concave when $\Sigma_n\succ0$.
\end{proof}
\begin{remark}
This is the sharpest correction to V1. It shows the pathology \emph{must} come from nonlinearity:
in the O'Brien et al.\ setting the ODE $\dot y=\lambda y$ is linear in $y$ but its numerical
solution is nonlinear in the \emph{parameter} $\lambda$. V1's ``sign-definite $\Rightarrow$ unimodal''
criterion is simply the linear case in disguise.
\end{remark}

\subsection{The mechanism: non-injectivity of the numerical map}\label{sec:rk4}
Consider $y=\varphi(\lambda)+\eta$, $\eta\sim N(0,\sigma^2)$, with the exact map
$\varphi(\lambda)=e^{\lambda T}$ and its RK4 surrogate with $n$ steps of size $h=T/n$,
\[
  \varphi^{\mathrm{RK4}}_n(\lambda)=R(\lambda h)^{n},\qquad
  R(z)=1+z+\tfrac{z^2}{2}+\tfrac{z^3}{6}+\tfrac{z^4}{24}.
\]
$e^{z}$ is strictly increasing, but $R$ is \emph{not}: $R'(z)=1+z+z^2/2+z^3/6$ has a root
$z^\star=-1.5960716\ldots\in(-2,-1)$ at which $R$ attains an interior minimum $R(z^\star)=0.2703948$.
Consequently $\lambda\mapsto R(\lambda h)^n$ is non-injective, with its fold at
\begin{equation}\label{eq:nstar}
  \lambda_{\mathrm{fold}}(n)=\frac{z^\star n}{T},
\end{equation}
while the exact map $e^{\lambda T}$ is strictly increasing. The fold is present for \emph{every} $n$
and simply moves out to $|\lambda|=|z^\star|n/T$; what decides whether the posterior is bimodal is
whether $\lambda_{\mathrm{fold}}(n)$ lies in the region carrying posterior mass, i.e.\ the prior's
effective support $[m_0-3\sigma_0,\,m_0+3\sigma_0]$. For the tables below ($m_0=-1$, $\sigma_0=0.5$,
$T=1$) this gives
\[
  |z^\star|\frac{n}{T}\lesssim |m_0|+3\sigma_0=2.5
  \quad\Longleftrightarrow\quad n\lesssim1.57,
\]
so observable-mass bimodality is confined to $n=1$; for larger $n$ the fold lies outside the
prior-supported region and the surviving secondary modes carry negligible mass (quantified below).

The trapezoidal rule is different: $T(z)=(1+z/2)/(1-z/2)$ has $T'(z)=1/(1-z/2)^2>0$, so its
surrogate is strictly increasing on the physical branch $z<2$ and cannot fold there. (It has a pole
at $z=2$ and, for even $n$, a second preimage in the far tail $\lambda h>2$, but with negligible
mass in the configurations considered.)

\begin{center}
\begin{tabular}{rlll}
\toprule
$n$ & exact & RK4 & trapezoidal \\
\midrule
$1$  & $1$ mode & \textbf{2 modes}: $-1.02$ and $-2.02$ (mass $0.11$) & $1$ mode \\
$2$  & $1$ mode & $2$ modes: $-1.00$ and $-4.71$ (mass $2\times10^{-13}$) & $1$ mode \\
$4$  & $1$ mode & $1$ mode & $1$ mode \\
$8$  & $1$ mode & $1$ mode & $1$ mode \\
$16$ & $1$ mode & $1$ mode & $1$ mode \\
$64$ & $1$ mode & $1$ mode & $1$ mode \\
\bottomrule
\end{tabular}
\end{center}

\noindent Table 1. All local maxima of the posterior for $\lambda$ ($T=1$, prior $N(-1,0.5^2)$,
$\sigma=0.05$, data $y=e^{-1}$), computed by the shipped driver with \emph{no} height threshold;
``mass'' is the relative \emph{peak height} of the secondary mode (the integrated relative mass at
$n=1$ is $0.094$). RK4 has a genuinely observable second
mode (mass $O(10^{-1})$) only at the coarsest level; by $n=2$ the secondary mode is already
exponentially suppressed, and it lies outside the prior-supported region for $n\ge4$.

\begin{remark}[Correction of V1's Experiment 5]
V1 reported that RK4 ``remains bimodal even at $n=64$''. Table~1 shows the opposite: an
observable-mass fold exists only while \eqref{eq:nstar} fails, so the RK4 secondary mode's mass drops
from $\approx0.11$ at $n=1$ to $\approx2\times10^{-13}$ at $n=2$ and is confined to the far prior tail
thereafter. Moreover V1's published code for this experiment does not integrate any ODE: it
hard-codes a bimodal \emph{mixture prior} $0.5N(-2,1)+0.5N(2,1)$ and ``refines'' by dividing the
observation noise by $\sqrt n$. That experiment therefore says nothing about discretisation error,
and V1's Table 8 has no reproducible source. We supply a genuine implementation in the driver.
\end{remark}

\begin{remark}[A correct diagnostic]
Equation \eqref{eq:nstar} gives practitioners a check that V1's ``sign of $F_{n+1}-F_n$'' does not:
compare the fold location $\lambda_{\mathrm{fold}}(n)=z^\star n/T$ with the prior's effective support
$[m_0\pm3\sigma_0]$. Bimodality is possible only while the fold lies in that range, i.e.\
$n\lesssim(|m_0|+3\sigma_0)T/|z^\star|$; it is not a fine-grid phenomenon.
\end{remark}

\section{Heavy-tailed priors: what is proved, and what is not}\label{sec:tails}
V1's Theorem 10.1 claimed a sharp rate $\rho(\alpha)(k+1)$ with $\rho(\alpha)=(\alpha-2)/(\alpha-1)$
for priors with polynomial tails $\sim\|x\|^{-\alpha}$. The proof is a truncation argument that does
not close (the ``$2p$ term cancels in the optimization'' step is asserted, never performed), the
matching lower bound is announced rather than derived, and the reference supporting it is mis-cited:
V1's \texttt{[38]} gives DOI \texttt{10.1214/24-AOS2356}, which resolves to \emph{Stankewitz et al.,
``Early stopping for $L^2$-boosting in high-dimensional linear models''}; the intended work is
Agapiou--Castillo, \emph{Heavy-tailed Bayesian nonparametric adaptation}, Ann.\ Statist.\
\textbf{52}(4) (2024), 1433--1459, DOI \texttt{10.1214/24-AOS2397}.

What \emph{is} true is the following sufficient condition, which is what V1's own Lemma 6.2 in fact
uses.

\begin{proposition}[Sufficient condition]\label{prop:suff}
If the uniform exponential-moment condition \eqref{eq:unif} holds, then the conclusion of
Theorem~\ref{thm:conv} holds with the full rate $O(n^{-(k+1)})$; no degradation occurs.
\end{proposition}
\begin{proof}
With $|\Delta\Phi_n|\le C(1+\|x\|^{2p})n^{-(k+1)}$ (Corollary~\ref{cor:incbound}) and Lemma
$|e^{-t}-1|\le|t|e^{|t|}$, one has $\E_{\mu_n}[|e^{-\Delta\Phi_n}-1|]
\le Cn^{-(k+1)}\E_{\mu_n}[(1+\|x\|^{2p})e^{C(1+\|x\|^{2p})n^{-(k+1)}}]$, which is
$O(n^{-(k+1)})$ uniformly provided \eqref{eq:unif} holds. The normalisation ratio
$Z_n/Z_{n+1}=1+O(n^{-(k+1)})$ follows the same way. Combining with the two-part split of V1's
proof of Theorem 6.1 gives the claim.
\end{proof}

\begin{remark}[The condition is sufficient, not necessary --- and V1's experiment shows no penalty]
Condition \eqref{eq:unif} fails for \emph{any} polynomial-tailed posterior (indeed for any
non-Gaussian tail), so it cannot by itself establish degradation; and it turns out not to force it
either. We reproduced V1's own Experiment 6 (which uses $F_n(\theta)=\theta(1-n^{-2})$, so $k=2$)
with priors Gaussian, $t_{10}$, $t_5$, $t_3$, $t_2$, computing the \emph{consecutive}-level
increment $|\E_{\mu_n}[\theta]-\E_{\mu_{n+1}}[\theta]|$:
\begin{center}
\begin{tabular}{lccccc}
\toprule
prior & Gaussian & $t_{10}$ & $t_5$ & $t_3$ & $t_2$ \\
\midrule
measured slope & $-2.97$ & $-2.97$ & $-2.97$ & $-2.97$ & $-2.97$ \\
V1's $\rho(\alpha)(k+1)$ & $-3.000$ & $-2.667$ & $-2.250$ & $-1.500$ & $\phantom{-}0.000$ \\
\bottomrule
\end{tabular}
\end{center}
The measured rate is $-(k+1)$ in every case; there is no degradation, and V1's Table 9 (which
reported $-2.61,-2.19,-1.47$ for the same priors) is not reproducible from its own model. What goes
wrong in V1 is the inference ``the sufficient condition fails $\Rightarrow$ the rate degrades''.
\end{remark}

\begin{conjecture}[Open problem]\label{conjecture}
Characterise the exact degradation of the telescoping rate under heavy-tailed priors: for which
forward-approximation families and which tail indices does $|\E_{\mu_n}[\varphi]-\E_{\mu_{n+1}}[\varphi]|$
fail to be $\Theta(n^{-(k+1)})$? We expect a genuine penalty to require the increment functional to
be unbounded in the tail (so that $\E_{\mu_n}[\|x\|^{2p}]$ is not uniformly bounded), which the
model above does not satisfy; but we have neither a proof of a positive rate in general nor a
counterexample to one. V1's formula $\rho(\alpha)=(\alpha-2)/(\alpha-1)$ remains unproved and is not
supported by numerics.
\end{conjecture}

\section{An extension, with corrected cost accounting}\label{sec:tkf}
V1 proposed a telescoping Kalman filter (TKF) refining the transition model $A_n$ in time and claimed
cost $O(T^{3/2})$ versus $O(T^2)$ for a fixed-model filter. That accounting is wrong: with the
schedule $n_t=\lceil\sqrt t\rceil$, the per-step cost is $n_t^2\asymp t$, so
$\sum_{t=1}^T n_t^2=\Theta(T^2)$ --- the same as $T\cdot n_T^2$. The driver reports
$\sum_{t\le10^4}n_t^2=5.07\times10^7$ against $T^{3/2}=10^6$. There is therefore no asymptotic
advantage; any gain is a constant-factor/early-transient effect, not a complexity improvement. We
retain the TKF as a heuristic with no asymptotic claim.

\section{Discussion}
The corrected picture is smaller but sound. Telescoping Bayesian inference gives, under clean
exponential moments, the rate $O(n^{-(k+1)})$ for posterior expectation differences
(Theorem~\ref{thm:conv}), the same rate for variance differences (Theorem~\ref{thm:var}), and the
standard multilevel cost with the correct threshold $\gamma=\beta+2$ (Theorem~\ref{thm:mlmc}). The
two ``non-ideal'' phenomena that motivated V1 behave quite differently from what V1 claimed:
discretisation-induced bimodality is impossible for linear forward operators
(Theorem~\ref{thm:obstruction}) and, when it does occur, is a coarse-resolution transient tied to
non-injectivity of the numerical map (\S\ref{sec:rk4}); and the heavy-tail ``penalty'' is, at least
in the model that was supposed to exhibit it, absent (\S\ref{sec:tails}). Both are genuine open
problems, now stated honestly.

\paragraph{Reproducibility.} Every number in this paper is regenerated by
\texttt{code/verify\_bt\_v2.py}, which asserts every claimed sign and threshold and exits nonzero on
failure; it reports \texttt{ALL ASSERTIONS PASS}. The manuscript, the driver and the two independent
vendor audit reports are archived at \url{https://github.com/jpbald93/bayesian-telescoping-v2}.
Unlike V1, we claim no reproduction beyond these artifacts, and we cite no experiment we did not run.

\section*{Change log relative to V1}
\begin{itemize}
\item \textbf{Retracted:} V1 Thm~9.3 (bimodality characterisation) --- invalid Hessian argument;
replaced by Thm~\ref{thm:obstruction} and \S\ref{sec:rk4}.
\item \textbf{Retracted:} V1 Thm~10.1 (heavy-tail penalty) --- unproved and unobserved; replaced by
Prop~\ref{prop:suff} and Conjecture~\ref{conjecture} (see \S\ref{sec:tails}).
\item \textbf{Corrected:} V1 Thm~8.4 thresholds $\beta+1\to\beta+2$ and exponent
$-2-(\beta+1-\gamma)/(k+1)\to-2-(\beta+2-\gamma)/k$.
\item \textbf{Corrected:} telescope tail lemma integral step (Lemma~\ref{lem:tail}).
\item \textbf{Corrected:} TKF cost $O(T^{3/2})\to\Theta(T^2)$.
\item \textbf{Corrected:} reference for heavy-tailed priors (DOI/venue/authors); the multimodal
reference in V1 cited arXiv:0909.2126, which is a different 2009 paper (the intended work is
O'Brien, Moores et al., arXiv:2502.11510, 2025).
\item \textbf{Removed:} all numerical tables not reproducible from a shipped script, and the claim
of a public reproduction repository.
\end{itemize}

\section*{Code and data availability}
The complete reproduction package --- \LaTeX{} source, compiled PDF, the verification driver
\texttt{code/verify\_bt\_v2.py} together with its recorded output, and the two independent vendor
audit reports --- is archived at \url{https://github.com/jpbald93/bayesian-telescoping-v2}. Running
\texttt{python3 verify\_bt\_v2.py} (NumPy + mpmath) regenerates every numeral in this paper and exits
nonzero if any asserted sign or threshold fails. No table, figure or numerical claim in this paper
depends on code or data outside that package.

\section*{Acknowledgments}
The author acknowledges the use of AI-assisted tools for editorial and programming support. All
mathematical statements, proofs and computations were verified against the shipped driver.

\begin{thebibliography}{9}
\bibitem{dashti} M.~Dashti and A.~M. Stuart, \emph{The Bayesian approach to inverse problems}, in
  Handbook of Uncertainty Quantification, Springer (2017), 311--428.
\bibitem{giles} M.~B. Giles, \emph{Multilevel Monte Carlo methods}, Acta Numer. \textbf{24} (2015),
  259--328.
\bibitem{dodwell} T.~J. Dodwell, C.~Ketelsen, R.~Scheichl, and A.~L. Teckentrup, \emph{Multilevel
  Markov chain Monte Carlo}, SIAM Rev. \textbf{61} (2019), 509--545.
\bibitem{jasra} A.~Jasra, K.~J.~H. Law, N.~Walton, and S.~Yang, \emph{Multi-index sequential Monte
  Carlo ratio estimators for Bayesian inverse problems}, Found.\ Comput.\ Math.\ \textbf{24} (2024),
  1249--1304.
\bibitem{latz} J.~Latz, \emph{Bayesian inverse problems are usually well-posed}, SIAM Rev.
  \textbf{65} (2023), 831--865.
\bibitem{obrien} T.~O'Brien, M.~T. Moores, et al., \emph{Here be dragons: bimodal posteriors arise
  from numerical integration error in longitudinal models}, arXiv:2502.11510 (2025).
\bibitem{ac} S.~Agapiou and I.~Castillo, \emph{Heavy-tailed Bayesian nonparametric adaptation},
  Ann.\ Statist.\ \textbf{52}(4) (2024), 1433--1459.
\end{thebibliography}

\appendix
\section{Verification driver}\label{app:code}
\texttt{code/verify\_bt\_v2.py} (Python; NumPy + mpmath;
\url{https://github.com/jpbald93/bayesian-telescoping-v2}) regenerates every numeral above and asserts:
the tail bound (Lemma~\ref{lem:tail}); the corrected MLMC regimes (Theorem~\ref{thm:mlmc}), including
the divergence at V1's threshold; monotonicity of $e^z$ and $T(z)$ and the interior minimum of $R$;
the RK4/trapezoidal mode counts of Table~1; and the heavy-tail slopes of \S\ref{sec:tails}. It exits
nonzero on any failed assertion and prints \texttt{ALL ASSERTIONS PASS}.

\end{document}
